\section{Discussion}

Three conditions decide whether learning from users works: the judge keeps failed attempts out of the lessons, the user can see the outcome so that this is possible, and lessons cover the skills evenly.

\begin{itemize}
\tightlist
\item
  \textbf{Precision matters more than volume.} A handful of correct lessons was enough to improve (5.6), while one failed attempt on night 1 could end a life. A judge should refuse to learn when in doubt. Why clean but sparse lessons still faded after two weeks (5.5) is open.
\item
  \textbf{Discernment is bounded by what the user can see.} The sense read customers correctly; the customers were wrong. No amount of discernment in reading a reaction can recover what the user did not know.
\item
  \textbf{Replay follows success.} Lessons come from what the model already does well, so replay narrows its skills; balancing by domain or task type is the obvious fix.
\item
  \textbf{Practical advice:} ask users for a verdict only when they can check the result; mix in verified outcomes where available; balance replay across skills.
\item
  \textbf{To write: future work.} A fix for forgetting (replay balanced by domain or task type, tested on $\tau^2$-bench); an explanation of the late fade (an \emph{oracle} limited to about 2 lessons a day); a test of the sense on real users' reactions; larger models.
\end{itemize}
